% GENERATED FILE. Edit canonical lessons, then run npm run generate:books.
% canonical-source-hash: fnv1a64:423133f0453914b2
% canonical-lessons: TE-C167-talanoppi, TE-C167-virecanalu, TE-C167-matra, TE-C167-cemata, TE-C167-vapu

\chapter{Headache, Diarrhoea, Tablet, Sweat, Swelling}
\label{ch:te-headache-diarrhoea-tablet-sweat-swelling}
\hlchaptermodality{167}

\begin{chapteropening}
\textbf{By the end of this chapter:} \emph{I can say a headache, diarrhoea, a tablet, sweat and a swelling.}

Complete the last lesson of chapter 167: I can say a headache, diarrhoea, a tablet, sweat and a swelling.
\end{chapteropening}

\section[\texorpdfstring{\te{తలనొప్పి} (talanoppi)}{తలనొప్పి (talanoppi)}]{\te{తలనొప్పి} (talanoppi) — a headache}

\label{lesson:TE-C167-talanoppi}



\begin{quote}
Before the new one: say the Telugu for to mutter, then the Telugu for to separate.
\end{quote}

\subsection*{You'll want to know: \te{తలనొప్పి}}
\textbf{\te{తలనొప్పి}} — \emph{talanoppi} — \textquotedblleft{}a headache\textquotedblright{}.

\begin{grammarlens}[title={What you've built}]
One more word for this chapter.
\end{grammarlens}

\subsection*{Your turn}
\begin{itemize}
\raggedright
  \item \emph{Say it:} \emph{talanoppi}
  \item \emph{Say it:} \emph{talanoppi}, once more
  \item \emph{Say it:} say \emph{talanoppi}
  \item \emph{Recall:} say the Telugu for to mutter, then the Telugu for to separate, then say \emph{talanoppi} again
\end{itemize}

\begin{tcolorbox}[breakable,colback=teal!4,colframe=teal!35!black,title={Before you move on}]
What does \te{తలనొప్పి} mean? (\textquotedblleft{}A headache\textquotedblright{}.) Say it once more.
\end{tcolorbox}
\section[\texorpdfstring{\te{విరేచనాలు} (virēcanālu)}{విరేచనాలు (virēcanālu)}]{\te{విరేచనాలు} (virēcanālu) — diarrhoea}

\label{lesson:TE-C167-virecanalu}



\begin{quote}
Before the new one: say the Telugu for to separate, then the Telugu for a headache.
\end{quote}

\subsection*{You'll want to know: \te{విరేచనాలు}}
\textbf{\te{విరేచనాలు}} — \emph{virēcanālu} — \textquotedblleft{}diarrhoea\textquotedblright{}.

\begin{grammarlens}[title={What you've built}]
One more word for this chapter.
\end{grammarlens}

\subsection*{Your turn}
\begin{itemize}
\raggedright
  \item \emph{Say it:} \emph{virēcanālu}
  \item \emph{Say it:} \emph{virēcanālu}, once more
  \item \emph{Say it:} say \emph{virēcanālu}
  \item \emph{Recall:} say the Telugu for to separate, then the Telugu for a headache, then say \emph{virēcanālu} again
\end{itemize}

\begin{tcolorbox}[breakable,colback=teal!4,colframe=teal!35!black,title={Before you move on}]
What does \te{విరేచనాలు} mean? (\textquotedblleft{}Diarrhoea\textquotedblright{}.) Say it once more.
\end{tcolorbox}
\section[\texorpdfstring{\te{మాత్ర} (mātra)}{మాత్ర (mātra)}]{\te{మాత్ర} (mātra) — a tablet, a pill}

\label{lesson:TE-C167-matra}



\begin{quote}
Before the new one: say the Telugu for a headache, then the Telugu for diarrhoea.
\end{quote}

\subsection*{You'll want to know: \te{మాత్ర}}
\textbf{\te{మాత్ర}} — \emph{mātra} — \textquotedblleft{}a tablet, a pill\textquotedblright{}.

\begin{grammarlens}[title={What you've built}]
One more word for this chapter.
\end{grammarlens}

\subsection*{Your turn}
\begin{itemize}
\raggedright
  \item \emph{Say it:} \emph{mātra}
  \item \emph{Say it:} \emph{mātra}, once more
  \item \emph{Say it:} say \emph{mātra}
  \item \emph{Recall:} say the Telugu for a headache, then the Telugu for diarrhoea, then say \emph{mātra} again
\end{itemize}

\begin{tcolorbox}[breakable,colback=teal!4,colframe=teal!35!black,title={Before you move on}]
What does \te{మాత్ర} mean? (\textquotedblleft{}A tablet, a pill\textquotedblright{}.) Say it once more.
\end{tcolorbox}
\section[\texorpdfstring{\te{చెమట} (cemaṭa)}{చెమట (cemaṭa)}]{\te{చెమట} (cemaṭa) — sweat}

\label{lesson:TE-C167-cemata}



\begin{quote}
Before the new one: say the Telugu for diarrhoea, then the Telugu for a tablet.
\end{quote}

\subsection*{You'll want to know: \te{చెమట}}
\textbf{\te{చెమట}} — \emph{cemaṭa} — \textquotedblleft{}sweat\textquotedblright{}.

\begin{grammarlens}[title={What you've built}]
One more word for this chapter.
\end{grammarlens}

\subsection*{Your turn}
\begin{itemize}
\raggedright
  \item \emph{Say it:} \emph{cemaṭa}
  \item \emph{Say it:} \emph{cemaṭa}, once more
  \item \emph{Say it:} say \emph{cemaṭa}
  \item \emph{Recall:} say the Telugu for diarrhoea, then the Telugu for a tablet, then say \emph{cemaṭa} again
\end{itemize}

\begin{tcolorbox}[breakable,colback=teal!4,colframe=teal!35!black,title={Before you move on}]
What does \te{చెమట} mean? (\textquotedblleft{}Sweat\textquotedblright{}.) Say it once more.
\end{tcolorbox}
\section[\texorpdfstring{\te{వాపు} (vāpu)}{వాపు (vāpu)}]{\te{వాపు} (vāpu) — a swelling}

\label{lesson:TE-C167-vapu}



\begin{quote}
Before the new one: say the Telugu for a tablet, then the Telugu for sweat.
\end{quote}

\subsection*{You'll want to know: \te{వాపు}}
\textbf{\te{వాపు}} — \emph{vāpu} — \textquotedblleft{}a swelling\textquotedblright{}.

\begin{grammarlens}[title={What you've built}]
One more word for this chapter.
\end{grammarlens}

\subsection*{Your turn}
\begin{itemize}
\raggedright
  \item \emph{Say it:} \emph{vāpu}
  \item \emph{Say it:} \emph{vāpu}, once more
  \item \emph{Say it:} say \emph{vāpu}
  \item \emph{Recall:} say the Telugu for a tablet, then the Telugu for sweat, then say \emph{vāpu} again
\end{itemize}

\begin{tcolorbox}[breakable,colback=teal!4,colframe=teal!35!black,title={Before you move on}]
What does \te{వాపు} mean? (\textquotedblleft{}A swelling\textquotedblright{}.) Say it once more.
\end{tcolorbox}
